\documentclass[11pt]{article}
\usepackage[T1]{fontenc}
\usepackage{lmodern,amsmath,amssymb,amsthm,mathtools}
\usepackage[margin=1.05in]{geometry}
\usepackage[hidelinks]{hyperref}
\newtheorem{theorem}{Theorem}[section]
\newtheorem{lemma}[theorem]{Lemma}
\newtheorem{proposition}[theorem]{Proposition}
\newtheorem{corollary}[theorem]{Corollary}
\newcommand{\R}{\mathbb R}\newcommand{\C}{\mathbb C}
\DeclareMathOperator{\Log}{Log}
\newcommand{\code}[1]{{\small\nolinkurl{#1}}}
\title{Eventual reality and simplicity of zeros\protect\\
for a growing sixth-order arithmetic transform}
\author{Project Stardust}
\date{Reviewed editorial draft, October 29, 2026}
\begin{document}
\maketitle
\begin{abstract}
We give a human proof of eventual reality and simplicity of
every zero of the cosine transform obtained by multiplying
the arithmetic theta kernel by $\exp(tu^6)$. The conclusion
holds on the entire complex plane, with one sufficiently
late time valid for all frequencies. The proof constructs
actual stationary points in a moving transition region and
its far complement, proves a uniform oscillatory Gaussian
estimate for the full kernel, and converts an analytic
cosine-pair approximation into exact real and simple zeros.
The sixth-order polynomial changes, physical-prefix bound,
and coverage by a common real stationary center are explicit.
The final theorem is matched to a frozen Lean development.
No general integer-order classification, optimal threshold
or Riemann-hypothesis consequence is asserted; novelty remains
unresolved.
\end{abstract}

\section{Statement and the actual integral}
For $u\geq0$ let
\begin{equation}\label{eq:kernel}
 K(u)=\sum_{n\geq1}
 (2\pi^2 n^4e^{9u}-3\pi n^2e^{5u})e^{-\pi n^2e^{4u}},
 \qquad
 F_t(z)=\int_0^\infty e^{tu^6}K(u)\cos(zu)\,du.
\end{equation}
The sign of the tilt is positive. Here and below $t$ is real
and $z$ is complex. A zero is simple when $F_t'(z)\ne0$.

\begin{theorem}\label{thm:main}
There is a finite $T$ such that, for every $t\geq T$ and
every $z\in\C$, $F_t(z)=0$ implies
\[
 \operatorname{Im}z=0,\qquad F_t'(z)\ne0.
\]
\end{theorem}

The time $T$ does not depend on $z$. In particular, a
fixed-frequency asymptotic cannot by itself prove this
statement. We use one escaping stationary center to cover
all frequencies. The proof follows the actual-saddle
organization of the quartic companion~\cite{quartic}, but
supplies the sixth-order derivatives and estimates needed
for the extension. It does not import the quartic arithmetic
certificate or its endpoint collision.

Each summand in \eqref{eq:kernel} is positive on the
half-line. Summation after splitting off half the
exponential gives
\begin{equation}\label{eq:realdecay}
 0<K(u)\leq C e^{9u-\frac\pi2e^{4u}},\qquad u\geq0.
\end{equation}
Indeed $n^2e^{4u}\geq(n^2+e^{4u})/2$ and
$\sum n^4e^{-\pi n^2/2}<\infty$. Thus every
$u^k e^{tu^6+B u}K(u)$ is integrable, for fixed $k,t,B$.
On compact parameter sets these are integrable dominating
functions for all spatial derivatives. Consequently $F_t$
is entire, even, conjugation compatible, and nonzero since
$F_t(0)>0$. All contour integrals below are ordinary
absolutely convergent integrals when their parameters are fixed.

\section{The theta strip and the sixth-order slope}
The kernel has an even holomorphic continuation to
\[
 \mathcal U=\{w:|\operatorname{Im}w|<\pi/8\}.
\]
Here is the precise kernel information we shall use. Write
$w=a+ib$, $E=e^{4w}$, $L=e^{4a}$ and $V=\operatorname{Re}E$.
For $a\geq0$, $V\geq1$, normal convergence of the series gives
\begin{equation}\label{eq:rho}
 K(w)=2\pi^2e^{9w-\pi E}(1+\rho(w)),\qquad
 |\rho^{(j)}(w)|\leq C_j\{e^{-4a}+(1+L)^j e^{-3\pi V}\},
 \quad 0\leq j\leq4.
\end{equation}
For completeness, after dividing by the leading exponential
the remaining factor is
\[
 1-\frac{3}{2\pi}e^{-4w}
 +\sum_{n\geq2}\left(n^4-\frac{3n^2}{2\pi}e^{-4w}\right)
       e^{-\pi(n^2-1)e^{4w}}.
\]
Each of its first four derivatives is bounded by a constant
times $(1+L)^j$ times a polynomial in $n$ and the same
exponential. For $n\geq2$, $n^2-1\geq3$; splitting off this
minimum and summing the remaining Gaussian proves
\eqref{eq:rho}.

The global strip bound needed at the physical endpoint is
\begin{equation}\label{eq:stripbound}
 |K(s+ib)|\leq C\cos(4b)^{-3}e^{-s},\qquad s\geq0.
\end{equation}
To check it, compare $\sum n^m e^{-\pi Vn^2}$, $m=2,4$,
with the Gaussian integral from zero, obtaining
$C_m V^{-(m+1)/2}$. Substitution of $V=e^{4s}\cos(4b)$
in the two terms of the kernel gives respectively
$C e^{-s}\cos(4b)^{-5/2}$ and
$C e^{-s}\cos(4b)^{-3/2}$, both bounded as stated.

Evenness follows from the Gaussian theta identity. If
$\theta(v)=\sum_{n\in\mathbb Z}e^{-\pi n^2v}$, its modular
identity is $\theta(v)=v^{-1/2}\theta(1/v)$, with the
right-half-plane square root; see~\cite{dlmf}. Put
$G(w)=e^w(\theta(e^{4w})-1)/2$. Then
$G(-w)=G(w)+\sinh w$, and direct differentiation gives
$(D^2-1)G=8K$. This proves evenness throughout $\mathcal U$.

For large $a$, under $V\geq c_0ta^5$, $t\geq1$, the
kernel is nonzero and its full logarithmic derivative is
\begin{equation}\label{eq:slope}
 \sigma_t(w)=6tw^5+\frac{K'(w)}{K(w)}
 =6tw^5+9-4\pi e^{4w}+\epsilon(w),\qquad
 |\epsilon^{(j)}(w)|\leq C e^{-4a}\quad(0\leq j\leq3).
\end{equation}
This follows by differentiating $\log(1+\rho)$ in
\eqref{eq:rho}. Indeed $L^4e^{-cV}\leq e^{16a-cta^5}$
is eventually smaller than $e^{-4a}$. The same bound holds
along $w+s$, $s\geq-1$, with $Ce^{-4a}$ on the right:
the exponential terms are bounded by
$\exp(16a+16s-cta^5e^{4s})$, whose maximum on this
half-line is at its left endpoint for large $a$.
All constants below are independent of $t$ and frequency,
unless a fixed parameter is indicated by a subscript.

\section{Uniform analysis at an actual saddle}
Fix $c_0=10^{-3}$. Call $(t,z,v)$ admissible when
\begin{equation}\label{eq:admissible}
 t\geq1,\quad v=a+ib\in\mathcal U,\quad a\text{ large},
 \quad y=\operatorname{Im}z\leq1,\quad
 \sigma_t(v)+iz=0,\quad V\geq c_0ta^5.
\end{equation}
These conditions will be proved in the constructions below;
no theta asymptotic is applied at an assumed model root.
Set
\[
 Y=\max(0,-y),\quad Q=ta^5+Y,\quad
 \beta=-\sigma_t'(v)/2,\quad W_z(w)=e^{tw^6+izw}K(w).
\]
Taking real parts in full stationarity gives
\begin{equation}\label{eq:stationaryreal}
 4\pi V=6t(a^5-10a^3b^2+5ab^4)+9-y+\operatorname{Re}\epsilon(v).
\end{equation}
Since $|b|<\pi/8$, this implies, uniformly,
\begin{equation}\label{eq:geometry}
 cQ\leq V\leq CQ,\quad Q\leq CL,\quad Q\geq a^5,
 \qquad cL\leq|\beta|\leq CL,\quad
 cQ\leq\operatorname{Re}\beta\leq CQ.
\end{equation}
In fact
\[
 \sigma_t'(v)=30tv^4-16\pi e^{4v}+\epsilon'(v),\qquad
 \sigma_t''(v)=120tv^3-64\pi e^{4v}+\epsilon''(v).
\]
The polynomial terms have size at most $Cta^4=o(Q)$,
whereas the exponential has modulus $L$ and real part $V$.
On any sufficiently small fixed horizontal interval the
second derivative is bounded by $CL$. On the entire
horizontal ray $s\geq-1$ one has
\begin{equation}\label{eq:curvature}
 \operatorname{Re}\sigma_t'(v+s)\leq-cQ.
\end{equation}
To see this, use $(a+s)^4\leq Ca^4e^{2s}$ for $s\geq-1$
and large $a$. Its contribution $Cta^4e^{2s}$ is absorbed
by $16\pi V e^{4s}$; their ratio is at most $C/a$ on
this ray. The error from \eqref{eq:slope} is negligible.

\begin{proposition}\label{prop:gaussian}
There is $\eta(a)\to0$ such that, uniformly at admissible saddles,
\begin{align}
 I_z(b)&:=\int_0^\infty W_z(s+ib)\,ds
   =W_z(v)\frac{\sqrt\pi}{\sqrt\beta}(1+\delta),
       &|\delta|&\leq\eta(a),\label{eq:gaussian}\\
 \int_0^\infty |W_z(s+ib)|\,ds
   &\leq C|W_z(v)|/\sqrt Q.\label{eq:mass}
\end{align}
The square root in \eqref{eq:gaussian} has positive real part.
\end{proposition}
\begin{proof}
On $s\geq-1$ put
$P(s)=\int_0^s(\sigma_t(v+r)+iz)\,dr$.
The exact scalar differential equation gives
$W_z(v+s)=W_z(v)e^{P(s)}$. Full stationarity and
\eqref{eq:curvature} imply
\[
 \operatorname{Re}P(s)\leq-cQs^2,\qquad
 P(s)=-\beta s^2+O(L|s|^3)\quad(|s|\leq\rho),
\]
where $\rho>0$ is fixed and small. In particular this
argument does not assume a fixed complex disk as wide as
the horizontal contour when $b$ is near the strip boundary.

Let $r=L^{-7/16}$. The integral on $[-r,r]$ differs
from the Gaussian integral there by $O(Lr^4)$, since
$Lr^3\to0$ and $|e^h-1|\leq e^{|h|}|h|$.
For $\operatorname{Re}\beta>0$,
\[
 \int_{\mathbb R}e^{-\beta s^2}\,ds=\sqrt\pi/\sqrt\beta,
 \qquad
 \left|\int_r^\infty e^{-\beta s^2}\,ds\right|
 +\left|\int_{-\infty}^{-r}e^{-\beta s^2}\,ds\right|
 \leq\frac{2}{|\beta|r}.
\]
The identity follows by differentiating in $\beta$ and
integrating $(se^{-\beta s^2})'$: the product of the integral
with $\sqrt\beta$ is constant and equals $\sqrt\pi$ on the
positive real axis. Integration by parts using
$(e^{-\beta s^2})'=-2\beta s e^{-\beta s^2}$ proves the
tail bound. These are oscillatory bounds in $|\beta|$,
not bounds that incorrectly compare its real part to its modulus.

On $r\leq|s|\leq\rho$, Taylor's theorem gives
$|P'(s)|\geq cL|s|$ and $|P''(s)|\leq CL$.
Integration of $(e^P)'/P'$ bounds these annuli by
$C/(Lr)$. The local remainder divided by the main Gaussian
is $O(L^{3/2}r^4)=O(L^{-1/4})$, and each normalized
Gaussian or annular tail is $O(L^{-1/16})$.
The remaining intervals $[-1,-\rho]$ and $[\rho,\infty)$
are bounded by $Ce^{-cQ}$, using the real-part estimate.
After normalization their bounds still tend to zero, since
$\sqrt L=e^{2a}$ and $Q\geq a^5$.

It remains to treat the physical prefix $0\leq s\leq a-1$.
Write $z=X+iy$. Since $\cos(4b)=V/L\geq cQ/L$,
\eqref{eq:stripbound} and
$\operatorname{Re}(s+ib)^6\leq s^6+C$ give
\begin{equation}\label{eq:prefixupper}
 |W_z(s+ib)|\leq C e^{12a}Q^{-3}
  \exp\{t(a-1)^6+Ct+Y(a-1)-Xb\}.
\end{equation}
On the other hand, \eqref{eq:rho} and \eqref{eq:stationaryreal}
give at the center
\begin{equation}\label{eq:centerlower}
 |W_z(v)|\geq
 \exp\{ta^6-\tfrac32ta^5-Cta^4+9a+Y(a-\tfrac14)-Xb-Ca\}.
\end{equation}
Indeed $\operatorname{Re}v^6=a^6-15a^4b^2+15a^2b^4-b^6$
and $\pi V=\tfrac32t\operatorname{Re}v^5+(9-y+\operatorname{Re}\epsilon)/4$.
The difference of the polynomial exponents is
\[
 t\{(a-1)^6-a^6+\tfrac32a^5+Ca^4+C\}
 =-\tfrac92ta^5+O(ta^4)\leq-cta^5.
\]
The $Y$ difference is $-3Y/4$. Thus the prefix integral,
divided by $|W_z(v)|$, is at most $Ca e^{-cQ+Ca}$.
It vanishes even after multiplication by $\sqrt L$.
This completes \eqref{eq:gaussian}. Absolute integration
of $e^{-cQs^2}$ and the same prefix bound proves
\eqref{eq:mass}. One may, for example, bound $\eta$ by a
constant times $L^{-1/4}+L^{-1/16}+(1+a)e^{-ca^5+Ca}$.
\end{proof}

\section{Contours, reflection, and exact local zeros}
Let $J_z=\int_0^\infty W_z(s)\,ds$. For real $h$ in the
theta strip, a rectangle gives
\[
 J_z=I_z(h)+i\int_0^hW_z(iy)\,dy.
\]
At its right edge the modulus tends to zero: the leading
term $-\pi e^{4R}\cos(4h)$ dominates $tR^6+O(R^5)$
and all frequency and kernel polynomial terms.
For $y=\operatorname{Im}z$ this yields the exact identities
\begin{align}
 F_t(z)&=\frac{I_z(h)+\overline{I_{\bar z}(k)}}2
          +\frac i2\int_k^h W_z(iq)\,dq,\label{eq:pair}\\
 F_t(z)&=\tfrac12\{I_z(h)+I_{-z}(-h)\}.\label{eq:mirror}
\end{align}
Evenness of $K$ and of the sixth power cancels the vertical
integrals in the second formula.

For an admissible saddle with $y\leq0$, the reflected
integrands satisfy the exact modulus relation
\[
 |W_{-z}(s-ib)|=e^{2ys}|W_z(s+ib)|.
\]
On $s\geq a-1$ use $e^{2ys}\leq e^{2y(a-1)}$ and
\eqref{eq:mass}. On the prefix use $e^{2ys}\leq1$ and
the prefix bound. With $A_z=W_z(v)\sqrt\pi/\sqrt\beta$,
\begin{equation}\label{eq:reflected}
 \left|I_{-z}(-b)/A_z\right|
 \leq C\sqrt{L/Q}\,e^{2y(a-1)}+o(1).
\end{equation}
In particular no such zero has $y\leq-1$ for large $a$:
$\sqrt L=e^{2a}$ cancels $e^{-2(a-1)}$, leaving
$C/\sqrt Q+o(1)$, whereas $I_z(b)/A_z\to1$.

\begin{lemma}\label{lem:cos}
Let $f,E,\Phi$ be holomorphic on $|z-x|\leq1$, $x$ real,
with conjugation compatibility and $E$ nonvanishing. Suppose,
for $a$ large and $D/a\to0$,
\[
 |\Phi'(z)-a|\leq D,\qquad
 |f(z)/E(z)-\cos\Phi(z)|\leq\varepsilon\cosh(\operatorname{Im}\Phi(z)),
 \qquad\varepsilon\to0.
\]
Then every zero with $|z-x|<1/2$ is real and simple.
\end{lemma}
\begin{proof}
At a zero $q$, write $\Phi(q)=u+iv$.
The identity $|\cos(u+iv)|^2=\cos^2u+\sinh^2v$
gives $\tanh|v|\leq\varepsilon$ and $|\cos u|\leq2\varepsilon$.
On the disk $|z-q|\leq\delta/a$, $\delta=10^{-3}$,
integration of the derivative bound gives
$|\Phi(z)-u|\leq2\delta+2\varepsilon$.
Take $\varepsilon$ sufficiently small. If $d=\sin u$, then
$|d|\geq1/2$ and $|\sin\Phi(z)-d|$ is bounded by a small
absolute constant. Cauchy's estimate on this disk bounds
the derivative of $f/E-\cos\Phi$ on its half-disk by
$C\varepsilon a/\delta$. Hence there
\[
 |(f/E)'(z)+ad|\leq a|d|/2.
\]
Also vertical integration of $\Phi'$ from the real axis
gives $|\operatorname{Im}\Phi(q)|\geq(a-D)|\operatorname{Im}q|$.
Thus $|q-\bar q|\leq C\varepsilon/a<\delta/(2a)$.
If $q$ were nonreal, $q$ and $\bar q$ would be distinct
zeros in the half-disk. Integration of the last derivative
bound along their segment gives
$|(f/E)(q)-(f/E)(\bar q)|\geq a|d||q-\bar q|/2>0$,
a contradiction. The same bound shows the derivative at
the real zero is nonzero. Multiplication by $E$ preserves simplicity.
\end{proof}

\begin{proposition}\label{prop:local}
If a real $x$ has an admissible saddle $v_0=a+ib$, then,
for large $a$, every zero $z$ with $|z-x|<1/2$ is real and simple.
\end{proposition}
\begin{proof}
The theta boundary has distance at least $cQ/L$ from $v_0$,
because $\cos(4b)\geq cQ/L$. This is much larger than
$1/L$. For $|z-x|\leq2$, solve $\sigma_t(v)+iz=0$ by
the contraction
\[
 v\longmapsto v-\frac{\sigma_t(v)+iz}{\sigma_t'(v_0)}
\]
on a disk of radius $C/L$ about $v_0$. Its derivative
has modulus $O(1/L)$, by the $CL$ bound on $\sigma_t''$,
and its center displacement is $O(1/L)$. It gives an
analytic branch with
$v(z)=v_0+O(1/L)$ and $v'(z)=-i/\sigma_t'(v(z))=O(1/L)$.
For $|z-x|\leq1$, the perturbed saddle remains admissible:
first its margin is at least $(c_0/2)ta^5$, then
\eqref{eq:stationaryreal} restores a fixed stronger margin
for large $a$. The quantities $a,L,Q$ remain comparable
to their central values.

The nonzero analytic function
$A(z)=W_z(v(z))\sqrt\pi/\sqrt{\beta(z)}$
has logarithmic derivative
\begin{equation}\label{eq:logder}
 \frac{A'(z)}{A(z)}=iv(z)-\frac{\beta'(z)}{2\beta(z)}
 =iv(z)+O(1/L).
\end{equation}
Full stationarity cancels the derivative of $v$ in the
exponential weight. Also $|\beta'|=O(1)$ because
$|\sigma_t''|=O(L)$ and $|v'|=O(1/L)$.
Choose $P=\log(A/A(x))$ with $P(x)=0$, put
$P^*(z)=\overline{P(\bar z)}$ and fix a real argument
$\alpha$ of $A(x)$. Define
\[
 E(z)=|A(x)|e^{(P+P^*)/2},\qquad
 \Phi(z)=\alpha+(P-P^*)/(2i).
\]
Both have the required conjugation property and
$\Phi'=a+O(1/L)$. Exactly,
\[
 \tfrac12(A(z)+\overline{A(\bar z)})=E(z)\cos\Phi(z),\qquad
 \tfrac12(|A(z)|+|A(\bar z)|)=|E(z)|\cosh(\operatorname{Im}\Phi(z)).
\]

Use \eqref{eq:pair} with heights $h=\operatorname{Im}v(z)$
and $k=\operatorname{Im}v(\bar z)$. Proposition~\ref{prop:gaussian}
controls its two horizontal integrals by these $A$ terms.
The heights differ by $O(|\operatorname{Im}z|/L)$, and
throughout their intervening interval
$\cos(4q)\geq cQ/L$. Equations \eqref{eq:stripbound}
and $\operatorname{Re}(iq)^6=-q^6\leq0$ give a bound
$\exp(Ca+t-Xh+C)$ for the residual vertical integral,
where $X=x$ may replace $\operatorname{Re}z$ at bounded
cost. At the central real saddle, stationarity implies
$|x|\leq CL$, so $|X(h-k)|=O(1)$.
By \eqref{eq:centerlower}, meanwhile,
$|A(z)|\geq\exp(ta^6-Cta^5-Ca-Xh)$.
The ratio tends uniformly to zero since $t\geq1$ and
$a\to\infty$. Thus
\[
 |F_t(z)/E(z)-\cos\Phi(z)|
 \leq o(1)\cosh(\operatorname{Im}\Phi(z)),\qquad |z-x|\leq1.
\]
Lemma~\ref{lem:cos} applies with $D=C/L$.
\end{proof}

\section{The moving transition region}
The positive continuous function $e^{tu^6}K(u)$ tends to
zero at infinity and attains a maximum on $[0,\infty)$.
Every such maximizer $u=u_t$ escapes to infinity as
$t\to\infty$. Indeed for a fixed $R$, the values on
$[0,R]$ are at most $C_R e^{tR^6}$, whereas the value
at $R+1$ is $K(R+1)e^{t(R+1)^6}$ and eventually larger.
At a positive maximizer differentiation gives
$\sigma_t(u)=0$. With $B=4\pi e^{4u}$,
\begin{equation}\label{eq:scale}
 B=6tu^5+9+\epsilon(u)\asymp tu^5.
\end{equation}
This is the center used in both frequency regions.

\begin{proposition}\label{prop:transition}
Fix $M>0$. For all sufficiently large $t$, every zero of
$F_t$ in $|z|\leq MB$ is real and simple.
\end{proposition}
\begin{proof}
It suffices first to take $\operatorname{Im}z\leq0$.
The anchor
$q=u+\tfrac14\Log(1+iz/B)$ has $q-u=O_M(1)$ and
a fixed positive distance, depending on $M$, from the
theta boundary. The exact fixed-point equation is
\begin{equation}\label{eq:transitionmap}
 w=u+\frac14\Log\left(1+\frac{iz}{B}
    +\frac{6t(w^5-u^5)+\epsilon(w)-\epsilon(u)}B\right).
\end{equation}
On the disk $|w-q|\leq C_M/u$, the logarithm argument
stays in a right half-plane: its leading real part is
at least one and its perturbation is $O_M(1/u)$.
The derivative of the perturbation is
$O_M(tu^4/B)=O_M(1/u)$. At $q$ the displacement is
$O_M(1/u)$; taking $C_M$ large and then $u$ large makes
the disk invariant and the map contractive. The root
is a full actual saddle with
\[
 \operatorname{Re}v=u+O_M(1),\qquad
 \operatorname{Re}e^{4v}\asymp_M B,\qquad
 e^{4\operatorname{Re}v}\asymp_M B.
\]
In particular the initial theta margin is positive; the
real stationarity equation strengthens it to
$\operatorname{Re}e^{4v}\geq cta^5$ with an absolute
$c>c_0$ for large $u$. Also $L/Q\leq C_M$.

For $y\leq-A_M/u$, the reflection bound
\eqref{eq:reflected} is at most $C_M e^{-A_M}+o(1)$,
since $a-1\geq u/2$. Choose $A_M$ large and then $u$
large to make this less than $1/4$, while the relative
error in \eqref{eq:gaussian} is less than $1/4$.
Equation \eqref{eq:mirror} excludes a zero. Conjugation
therefore confines every zero in this region to
$|\operatorname{Im}z|<A_M/u$.

For $x=\operatorname{Re}z$, $|x|\leq MB$, the same
construction supplies an admissible real-frequency saddle.
Once $u$ is large enough that $A_M/u<1/2$,
Proposition~\ref{prop:local} proves the zero real and
simple. Every estimate is uniform in this bounded-ratio region.
\end{proof}

\section{Actual far saddles and whole-plane coverage}
\begin{lemma}\label{lem:far}
There is an absolute $C_0$ such that, if $t\geq1$,
$X=|z|$, $\operatorname{Im}z\leq0$,
$\alpha=\tfrac14\log(X/(4\pi))$ is sufficiently large and
\begin{equation}\label{eq:farcondition}
 X\geq C_0t(\alpha+1)^5,
\end{equation}
then $z$ has an admissible actual saddle with real part
$\alpha+O(1)$ and $e^{4a}\asymp X$.
\end{lemma}
\begin{proof}
Start with the leading model, not the actual kernel. The
anchor $q=\tfrac14\Log(iz/(4\pi))$ has real part
$\alpha$ and imaginary part in $[-\pi/8,\pi/8]$,
since $\operatorname{Re}(iz)\geq0$. Solve
\[
 w=\tfrac14\Log\{(iz+6tw^5+9)/(4\pi)\}
\]
on the disk $|w-q|\leq Ct(\alpha+1)^5/X$.
Choose $C_0$ so large that this radius is less than $1/100$.
On the disk $|\operatorname{Im}w|<1/2$, and
$\operatorname{Re}w^5\geq(\operatorname{Re}w)^5/2$
for large $\alpha$. Hence the logarithm argument lies
strictly in the right half-plane. Its modulus is comparable
to $X$, and the map's derivative is
$O(t(\alpha+1)^4/X)=O(1/\alpha)$.
The center displacement is bounded by $Ct(\alpha+1)^5/X$,
using $|\Log(1+h)|\leq2|h|$ for $|h|\leq1/2$.
Increasing the fixed radius constant gives an invariant
contraction disk. Its fixed point solves
\[
 6tw^5+9-4\pi e^{4w}=-iz.
\]
If $w=a+ib$ and $Y=-\operatorname{Im}z$, its real part is
\[
 4\pi\operatorname{Re}e^{4w}
   =6t(a^5-10a^3b^2+5ab^4)+9+Y\geq cta^5+Y.
\]
Since $|b|<1/2$ and this quantity is positive, in fact
$|b|<\pi/8$. Also $e^{4a}\asymp X$. Thus the model root
lies strictly inside the theta strip with the required
positive margin; boundary anchors cause no missing case.

Now restore the kernel error. The model derivative at its
root is
\[
 30tw^4-16\pi e^{4w}
 =-4iz+t(30w^4-24w^5)-36.
\]
Its modulus is between $cX$ and $CX$, by
\eqref{eq:farcondition} with $C_0$ large.
On a disk of radius $1/X$, the variation is $O(1)$,
while $\epsilon$ and its first derivative are $O(e^{-4a})$.
Subtracting the full slope divided by the central model
derivative gives a contraction: its center displacement
is $O(e^{-4a}/X)$, and its derivative is $O(1/X)$.
The disk remains inside the strip because its boundary
distance is at least $c(ta^5+Y)/X\gg1/X$.
The resulting actual root preserves the positive margin;
\eqref{eq:stationaryreal} restores the fixed admissibility
constant. This proves the lemma.
\end{proof}

\begin{proposition}\label{prop:farzeros}
For an absolute sufficiently large $C_1$, the conditions
$t\geq1$, $\alpha$ sufficiently large and
$|z|\geq C_1t(\alpha+1)^5$ imply that every zero at $z$
is real and simple. For each finite $T$, all nonreal zeros
at $1\leq t\leq T$ lie in one finite disk.
\end{proposition}
\begin{proof}
Take $C_1\geq2C_0$. For $\operatorname{Im}z\leq-1$,
Lemma~\ref{lem:far} and \eqref{eq:reflected} exclude a
zero. Conjugation excludes $\operatorname{Im}z\geq1$.
If $|\operatorname{Im}z|<1$, let $x=\operatorname{Re}z$.
For large $|z|$, $|x|\geq|z|/2$, and
$\alpha_x=\tfrac14\log(|x|/(4\pi))\leq\alpha$.
Thus $|x|\geq C_0t(\alpha_x+1)^5$. Lemma~\ref{lem:far}
supplies a saddle at $x$. The local pair estimates in its
proof hold on $|q-x|\leq1$. At any proposed zero there,
the inequality $|\cos\Phi|\leq\varepsilon\cosh\operatorname{Im}\Phi$
gives $|\operatorname{Im}\Phi|\leq2\varepsilon$.
Vertical integration of $\Phi'=a+O(1/L)$ yields
$|\operatorname{Im}z|\leq2\varepsilon/(a-O(1/L))<1/2$.
Proposition~\ref{prop:local} now applies.
For $1\leq t\leq T$, the displayed power condition holds
uniformly for all sufficiently large $|z|$, since
$(\log X)^5/X\to0$. This proves the last assertion.
\end{proof}

\begin{proof}[Proof of Theorem~\ref{thm:main}]
Use a real maximizing center $u$ and $B=4\pi e^{4u}$ from
\eqref{eq:scale}. If $X=|z|\geq MB$, write
$\alpha=u+d$, so $d\geq\tfrac14\log M$. Then
\[
 \frac{t(\alpha+1)^5}{X}
 \leq C e^{-4d}\left(1+\frac{d+1}{u}\right)^5
 \leq C(2+d)^5e^{-4d},\qquad u\geq1.
\]
Choose one fixed $M>1$ so large that the last expression
is at most $1/C_1$ for every $d\geq\tfrac14\log M$.
For sufficiently late $t$, the center $u$ and hence
$\alpha$ are large, so Proposition~\ref{prop:farzeros}
settles all $X\geq MB$. For this same fixed $M$, enlarge
the time once more so Proposition~\ref{prop:transition}
settles all $X\leq MB$. These regions cover the entire
complex plane, using the same $u$ and $B$. This proves
the common-time conclusion.
\end{proof}

\section{Exact formal correspondence and publication scope}
The frozen mathematical baseline is
\begin{center}\small\ttfamily
bab688bb6140cd6316c80680095a914ea0d346db.
\end{center}
The exact final declaration is
\begin{center}\small
\code{RiemannSixthTransitionSimpleZeros.}\\
\code{eventually_all_zeros_real_simple}.
\end{center}
Its namespace prefix is
\begin{center}\code{OperatorDerivativeRewrite.Foundation}.\end{center}
The formal transform is
\begin{center}\code{RiemannEvenFlow.flow 3 t z}:\end{center}
the index three means a sixth power, not a third power.
The proof combines the bounded-ratio theorem with
\begin{center}\code{RiemannSixthFarZeroControl.}\\
\code{uniform_actual_far_zeros_real_simple},\end{center}
both at a single stationary center. The theorem map records
exact source lines and hashes. The human proof consolidates
explicit formal constants into uniform inequalities;
intermediate regrouped lemmas are not claimed as new named
formal statements.

De Bruijn~\cite{debruijn} and Newman~\cite{newman} provide
the primary historical comparison for real-zero Fourier
transforms. Newman's quadratic propagation theorem does
not supply the sixth-order conclusion. Targeted priority
review remains necessary; novelty is unresolved.

The theorem gives an existential late time. It gives no
numerical optimal threshold, no nonreal-zero certificate at
a specified early sixth-order time, no assertion that the
set of all reality times is an upper ray, and no threshold
collision. The argument is presented for the verified
sixth-order transform only. It does not certify every
integer even order, a real-order interpolation, or a
Riemann-hypothesis implication.

\begin{thebibliography}{9}
\bibitem{quartic} Project Stardust,
\emph{An attained eventual reality threshold for a growing
quartic arithmetic transform}, reviewed editorial draft,
October 29, 2026. Companion manuscript in
\code{manuscripts/quartic-threshold/}.
\bibitem{debruijn} N. G. de Bruijn,
The roots of trigonometric integrals,
\emph{Duke Math. J.} \textbf{17} (1950), 197--226.
\href{https://doi.org/10.1215/S0012-7094-50-01720-0}
{doi:10.1215/S0012-7094-50-01720-0}.
\bibitem{newman} C. M. Newman,
Fourier transforms with only real zeros,
\emph{Proc. Amer. Math. Soc.} \textbf{61} (1976), 245--251.
\url{https://sites.math.northwestern.edu/~auffing/papers/Newman.pdf}.
\bibitem{dlmf} NIST Digital Library of Mathematical Functions,
equation 20.7.32, modular theta identity.
\url{https://dlmf.nist.gov/20.7.E32}.
\end{thebibliography}
\end{document}
